% !TEX program = xelatex
% ============================================================================
%  POSN Computer Qualification — Real Numbers Slide Deck
%  Topic: จำนวนจริง (Real Numbers)
% ============================================================================

\documentclass[11pt, aspectratio=169]{beamer}

% --- Load shared preamble ---
% ============================================================================
%  POSN Computer Qualification — Shared Beamer Preamble
%  Used by all topic slide decks. Compile with XeLaTeX.
% ============================================================================

% ---------- Thai language & font ----------
\usepackage{iftex}
\usepackage{fontspec}

% XeTeX-specific: enable Thai line breaking
\ifXeTeX
  \XeTeXlinebreaklocale{th}
\fi
% Noto Sans Thai — body text (clean modern sans-serif, handles all Thai)
\setmainfont{Noto Sans Thai}[
  UprightFont = Noto Sans Thai Light,
  BoldFont    = Noto Sans Thai Medium,
  Scale       = 0.9
]
\setsansfont{Noto Sans Thai}[
  UprightFont = Noto Sans Thai Light,
  BoldFont    = Noto Sans Thai Medium,
  Scale       = 0.9
]

% Noto Sans Thai — clean modern sans-serif for structural elements
\newfontfamily\headingfont{Noto Sans Thai}[
  UprightFont = Noto Sans Thai Medium,
  BoldFont    = Noto Sans Thai Bold,
  Scale       = 0.9
]

% --- Heading font for structural elements ---
\setbeamerfont{title}{family=\headingfont, series=\bfseries, size=\huge}
\setbeamerfont{subtitle}{family=\headingfont, series=\mdseries}
\setbeamerfont{author}{family=\headingfont, series=\mdseries}

% --- Custom Title Page ---
\defbeamertemplate{title page}{custom}[1][]{%
  \centering
  \vspace*{2cm}
  % Title — in white
  {\usebeamerfont{title}\color{white}\inserttitle\par}
  \vspace{0.6cm}
  % Subtitle — in white
  {\usebeamerfont{subtitle}\color{white}\insertsubtitle\par}
  \vspace{2.5cm}
  % Author — blue filled rounded box, white text
  \begin{center}
    \begin{tcolorbox}[arc=3mm, colback=boxBlue, colframe=boxBlue!80,
                      width=0.42\textwidth, halign=center,
                      left=1.5em, right=1.5em, top=0.8em, bottom=0.8em,
                      boxrule=1.5pt, coltext=white]
      \usebeamerfont{author}\insertauthor
    \end{tcolorbox}
  \end{center}
}
\setbeamertemplate{title page}[custom]
\setbeamerfont{frametitle}{family=\headingfont, series=\bfseries}
\setbeamerfont{section in toc}{family=\headingfont}
\setbeamerfont{page number in head/foot}{family=\headingfont}
\setbeamerfont{section title}{family=\headingfont, series=\bfseries}
\setbeamerfont{subsection title}{family=\headingfont}

% Monospace font (for code/verbatim)
\setmonofont{Latin Modern Mono}[Scale=0.9]

% ---------- Math ----------
\usepackage{amsmath, amssymb}

% ---------- Graphics ----------
\usepackage{graphicx}
\usepackage{tikz}

% ---------- String parsing (for section headers) ----------
\usepackage{xstring}

% ---------- Colored boxes ----------
\usepackage[skins, breakable, xparse]{tcolorbox}
% Note: we do NOT load the 'most' option — it predefines example/solution environments
% that would conflict with our custom boxes.

% --- Color definitions ---
\definecolor{boxBlue}{RGB}{41, 128, 185}
\definecolor{boxGreen}{RGB}{39, 174, 96}
\definecolor{boxOrange}{RGB}{230, 126, 34}
\definecolor{boxPurple}{RGB}{142, 68, 173}
\definecolor{boxBlueBg}{RGB}{235, 245, 255}
\definecolor{boxGreenBg}{RGB}{235, 255, 240}
\definecolor{boxOrangeBg}{RGB}{255, 245, 235}
\definecolor{boxPurpleBg}{RGB}{248, 240, 255}

% --- Explanation box (blue) — for definitions, concepts, notes ---
\newtcolorbox{explanation}[1][]{
  enhanced,
  breakable,
  arc         = 2mm,
  colback     = boxBlueBg,
  colframe    = boxBlue!50,
  title       = #1,
  fonttitle   = \bfseries\color{white},
  coltitle    = white,
  colbacktitle= boxBlue,
  attach boxed title to top left = {yshift=-2mm, xshift=2mm},
  boxed title style = {arc=2mm, size=small},
  before skip = 8pt,
  after skip  = 8pt
}

% --- Example box (green) — for worked examples ---
\newtcolorbox{exambox}[1][]{
  enhanced,
  breakable,
  arc         = 2mm,
  colback     = boxGreenBg,
  colframe    = boxGreen!50,
  title       = #1,
  fonttitle   = \bfseries\color{white},
  coltitle    = white,
  colbacktitle= boxGreen,
  attach boxed title to top left = {yshift=-2mm, xshift=2mm},
  boxed title style = {arc=2mm, size=small},
  before skip = 8pt,
  after skip  = 8pt
}

% --- Solution box (orange) — for solutions and answers ---
\newtcolorbox{solbox}[1][]{
  enhanced,
  breakable,
  arc         = 2mm,
  colback     = boxOrangeBg,
  colframe    = boxOrange!50,
  title       = #1,
  fonttitle   = \bfseries\color{white},
  coltitle    = white,
  colbacktitle= boxOrange,
  attach boxed title to top left = {yshift=-2mm, xshift=2mm},
  boxed title style = {arc=2mm, size=small},
  before skip = 8pt,
  after skip  = 8pt
}

% --- Intuition box (purple) — for plain-language explanations, analogies ---
\newtcolorbox{intuition}[1][]{
  enhanced,
  breakable,
  arc         = 2mm,
  colback     = boxPurpleBg,
  colframe    = boxPurple!50,
  title       = #1,
  fonttitle   = \bfseries\color{white},
  coltitle    = white,
  colbacktitle= boxPurple,
  attach boxed title to top left = {yshift=-2mm, xshift=2mm},
  boxed title style = {arc=2mm, size=small},
  before skip = 8pt,
  after skip  = 8pt
}

% ---------- Beamer theme ----------
\usetheme{Madrid}
\usecolortheme{default}

% --- Custom beamer colors (blue accent matching explanation boxes) ---
\setbeamercolor{structure}{fg=boxBlue}
\setbeamercolor{palette primary}{fg=white, bg=boxBlue}
\setbeamercolor{palette secondary}{fg=white, bg=boxBlue!80}
\setbeamercolor{palette tertiary}{fg=white, bg=boxBlue!60}
\setbeamercolor{block title}{fg=white, bg=boxBlue}
\setbeamercolor{block body}{fg=black, bg=boxBlueBg}

% Remove navigation symbols for clean look
\setbeamertemplate{navigation symbols}{}

% Note: Frame title prefix removed — \addtobeamertemplate conflicts with beamer internals

% --- Outline (Table of Contents) styling ---
\setbeamerfont{section in toc}{family=\headingfont, series=\mdseries}
\setbeamerfont{subsection in toc}{family=\headingfont, series=\mdseries}
\setbeamertemplate{section in toc}{%
  \leavevmode\leftskip=1.5em
  \textcolor{boxBlue}{\textbf{\large\inserttocsectionnumber.}}%
  \quad\inserttocsection\par
  \vspace{0.2em}
}
\setbeamertemplate{subsection in toc}{%
  \leavevmode\leftskip=4em
  \textcolor{boxBlue!60}{\small$\bullet$}%
  \quad\inserttocsubsection\par
  \vspace{0.1em}
}

% Section header slides — Thai in white, English in smaller light blue below
\AtBeginSection{
  \begin{frame}[plain]{}
    \vspace*{2cm}
    \begin{beamercolorbox}[sep=12pt, center, rounded=true, shadow=true]{palette primary}
      \IfSubStr{\secname}{(}{%
        \StrBefore{\secname}{(}[\sectionthai]%
        \StrBetween[1,1]{\secname}{(}{)}[\sectioneng]%
        \begin{tabular}{@{}c@{}}
          {\usebeamerfont{title}\sectionthai} \\[0.2em]
          {\usebeamerfont{subtitle}\color{boxBlue!60}\sectioneng}
        \end{tabular}%
      }{%
        {\usebeamerfont{title}\secname\par}%
      }
    \end{beamercolorbox}
    \vfill
  \end{frame}
}

% Subsection header slides — smaller, centered
\AtBeginSubsection{
  \begin{frame}[plain]{}
    \vfill
    \centering
    {\usebeamerfont{section title}\insertsubsectionhead\par}
    \vfill
  \end{frame}
}

% Minimal footer: page number only, right-aligned
\setbeamertemplate{footline}{
  \hfill\usebeamerfont{page number in head/foot}\insertframenumber{} / \inserttotalframenumber\hspace*{1em}\vspace*{0.5ex}
}

% ---------- Spacing & readability ----------
\linespread{1.2}

% ---------- Useful commands ----------

% Highlight important text in blue
\newcommand{\highlight}[1]{\textcolor{boxBlue}{\textbf{#1}}}

% Math set notation helpers
\newcommand{\R}{\mathbb{R}}
\newcommand{\N}{\mathbb{N}}
\newcommand{\Z}{\mathbb{Z}}
\newcommand{\Q}{\mathbb{Q}}

% ============================================================================
%  End of preamble
% ============================================================================


% --- Document info ---
\title[จำนวนจริง]{\textcolor{white}{จำนวนจริง} \textcolor{boxBlue!60}{(Real Numbers)}}
\subtitle{เศษส่วน • ทศนิยม • อัตราส่วน • ร้อยละ • สมบัติจำนวนจริง • ค่าสัมบูรณ์ }
\author[None W. (Yuan)]{None W. (Yuan)}
\date{}

\begin{document}

% ===========================================================================
%  TITLE SLIDE
% ===========================================================================
\begin{frame}[plain]
  \titlepage
\end{frame}

% ===========================================================================
%  OUTLINE
% ===========================================================================
\begin{frame}{สารบัญ (Outline)}
  \tableofcontents
\end{frame}

% ===========================================================================
%  ก่อนเรียน (Prerequisites)
% ===========================================================================
\begin{frame}{ก่อนเรียน (Prerequisites)}
  \begin{explanation}[สิ่งที่ควรรู้ก่อนเรียน]
    \begin{itemize}
      \item \textbf{การดำเนินการพื้นฐาน} — บวก ลบ คูณ หาร
      \item \textbf{เซต} — ความเข้าใจเกี่ยวกับสมาชิกของเซต
    \end{itemize}
  \end{explanation}

  \begin{intuition}[จำนวนจริงคืออะไร?]
    จำนวนจริงคือจำนวนทุกชนิดที่เราพบในชีวิตประจำวัน — จำนวนนับ, จำนวนเต็ม,
    เศษส่วน, ทศนิยม, และจำนวนอตรรกยะอย่าง $\sqrt{2}, \pi$
  \end{intuition}
\end{frame}

% ===========================================================================
%  SECTION 1: ระบบจำนวนจริง
% ===========================================================================
\section{ระบบจำนวนจริง (Number System)}

\begin{frame}{ระบบจำนวนจริง}
  \begin{explanation}[แผนภาพระบบจำนวน]
    \textbf{จำนวนนับ} $\N = \{1, 2, 3, \ldots\}$

    \textbf{จำนวนเต็ม} $\Z = \{\ldots, -2, -1, 0, 1, 2, \ldots\}$

    \textbf{จำนวนตรรกยะ} $\Q = \{\frac{a}{b} \mid a, b \in \Z, b \neq 0\}$

    \textbf{จำนวนจริง} $\R$ — รวมทั้งตรรกยะและอตรรกยะ

    \medskip
    \highlight{$\N \subset \Z \subset \Q \subset \R$}
  \end{explanation}
\end{frame}

\begin{frame}{จำนวนตรรกยะ vs อตรรกยะ}
  \begin{explanation}[จำนวนตรรกยะ (Rational)]
    เขียนในรูป \highlight{$\frac{a}{b}$} โดย $a, b \in \Z$ และ $b \neq 0$
    \begin{itemize}
      \item ทศนิยมรู้จบ: $0.5 = \frac{1}{2}$, $0.125 = \frac{1}{8}$
      \item ทศนิยมซํ้า: $0.\overline{3} = \frac{1}{3}$, $0.\overline{142857} = \frac{1}{7}$
    \end{itemize}
  \end{explanation}

  \begin{explanation}[จำนวนอตรรกยะ (Irrational)]
    \textbf{ไม่สามารถ}เขียนในรูป $\frac{a}{b}$ ได้
    \begin{itemize}
      \item $\sqrt{2}, \sqrt{3}, \sqrt{5}, \ldots$
      \item $\pi \approx 3.14159\ldots$
      \item $e \approx 2.71828\ldots$
    \end{itemize}
  \end{explanation}
\end{frame}

\begin{frame}{กรณีพิเศษ: ทศนิยมซํ้า $0.\overline{1}$}
  \begin{exambox}[ตัวอย่าง: เขียน $0.111\ldots$ เป็นเศษส่วน]
    ให้ $x = 0.111\ldots$
    \begin{align*}
      10x &= 1.111\ldots \\
      x   &= 0.111\ldots \\
      \hline
      9x  &= 1 \quad \Rightarrow \quad x = \tfrac{1}{9}
    \end{align*}
  \end{exambox}

  \begin{intuition}[เทคนิคการแปลงทศนิยมซํ้า]
    \begin{itemize}
      \item ถ้าซํ้า 1 หลัก $\rightarrow$ คูณด้วย $10^1 = 10$ แล้วลบ
      \item ถ้าซํ้า 2 หลัก $\rightarrow$ คูณด้วย $10^2 = 100$ แล้วลบ
      \item ถ้าซํ้า $n$ หลัก $\rightarrow$ คูณด้วย $10^{\,n}$ แล้วลบ
    \end{itemize}
  \end{intuition}
\end{frame}

\begin{frame}{กรณีพิเศษ: อตรรกยะยกกำลังอาจเป็นตรรกยะ}
  \begin{exambox}[ตัวอย่างสุดคลาสสิก: $(\sqrt{2}^{\sqrt{2}})^{\sqrt{2}}$]
    พิจารณา $(\sqrt{2}^{\sqrt{2}})^{\sqrt{2}}$

    \medskip
    ใช้สมบัติ $(a^b)^c = a^{bc}$:
    \[
      (\sqrt{2}^{\sqrt{2}})^{\sqrt{2}} = \sqrt{2}^{\sqrt{2} \cdot \sqrt{2}} = \sqrt{2}^{\,2} = 2
    \]

    \medskip
    $\sqrt{2}$ เป็นอตรรกยะ, $\sqrt{2}^{\sqrt{2}}$ อาจเป็นอตรรกยะหรือไม่ก็ได้

    แต่ผลลัพธ์ $2$ เป็นจำนวนตรรกยะแน่นอน!

    \medskip
    \highlight{บทเรียน:} อตรรกยะยกกำลังอตรรกยะ สามารถให้ผลเป็นตรรกยะได้
  \end{exambox}
\end{frame}

% ===========================================================================
%  SECTION 2: เศษส่วน
% ===========================================================================
\section{เศษส่วน (Fractions)}

\begin{frame}{ประเภทของเศษส่วน}
  \begin{columns}[T]
    \column{0.48\textwidth}
      \textbf{เศษส่วนแท้ (Proper Fraction)}
      \begin{itemize}
        \item ตัวเศษ < ตัวส่วน
        \item ค่าน้อยกว่า 1
        \item เช่น $\frac{2}{3}, \frac{5}{8}, \frac{1}{4}$
      \end{itemize}

      \medskip
      \textbf{เศษส่วนเกิน (Improper Fraction)}
      \begin{itemize}
        \item ตัวเศษ > ตัวส่วน
        \item ค่ามากกว่า 1
        \item เช่น $\frac{7}{4}, \frac{5}{3}, \frac{11}{5}$
      \end{itemize}

    \column{0.48\textwidth}
      \textbf{จำนวนคละ (Mixed Number)}
      \begin{itemize}
        \item จำนวนเต็ม + เศษส่วนแท้
        \item เช่น $2\frac{1}{3}, 1\frac{3}{4}$
      \end{itemize}

      \medskip
      \textbf{การแปลง:}
      \[
        2\tfrac{1}{3} = \frac{2 \times 3 + 1}{3} = \frac{7}{3}
      \]
      \[
        \frac{11}{4} = 2\tfrac{3}{4} \quad (11 \div 4 = 2 \text{ เศษ } 3)
      \]
  \end{columns}
\end{frame}

\begin{frame}{การดำเนินการของเศษส่วน — บวกและลบ}
  \begin{exambox}[การบวก/ลบ — ต้องทำส่วนให้เท่ากันก่อน]
    \[
      \frac{1}{3} + \frac{1}{4} = \frac{4}{12} + \frac{3}{12} = \frac{7}{12}
    \]
    \[
      \frac{5}{6} - \frac{1}{2} = \frac{5}{6} - \frac{3}{6} = \frac{2}{6} = \frac{1}{3}
    \]
  \end{exambox}

  \begin{intuition}[ขั้นตอน]
    \begin{enumerate}
      \item หา ค.ร.น. ของตัวส่วน
      \item ทำตัวส่วนให้เท่ากัน
      \item บวก/ลบ ตัวเศษ (ตัวส่วนคงเดิม)
      \item ทำเป็นเศษส่วนอย่างตํ่า
    \end{enumerate}
  \end{intuition}
\end{frame}

\begin{frame}{การดำเนินการของเศษส่วน — คูณและหาร}
  \begin{exambox}[การคูณ/หาร]
    \[
      \frac{2}{3} \times \frac{4}{5} = \frac{2 \times 4}{3 \times 5} = \frac{8}{15}
    \]
    \[
      \frac{2}{3} \div \frac{4}{5} = \frac{2}{3} \times \frac{5}{4} = \frac{10}{12} = \frac{5}{6}
    \]
  \end{exambox}

  \begin{intuition}[เคล็ดลับ]
    \begin{itemize}
      \item \textbf{คูณ:} เศษ $\times$ เศษ, ส่วน $\times$ ส่วน
      \item \textbf{หาร:} กลับเศษเป็นส่วนของตัวหาร แล้วเปลี่ยนเป็นคูณ
    \end{itemize}
  \end{intuition}
\end{frame}

\begin{frame}{การทำให้เศษส่วนเป็นอย่างตํ่า}
  \begin{explanation}[วิธีทำ]
    \textbf{หารทั้งเศษและส่วนด้วย ห.ร.ม.} (GCD — ตัวหารร่วมมาก)
  \end{explanation}

  \begin{exambox}[ตัวอย่าง]
    \begin{align*}
      \frac{24}{36} &= \frac{24 \div 12}{36 \div 12} = \frac{2}{3} \quad (\text{ห.ร.ม. ของ 24, 36 คือ 12}) \\[1em]
      \frac{45}{60} &= \frac{45 \div 15}{60 \div 15} = \frac{3}{4} \quad (\text{ห.ร.ม. ของ 45, 60 คือ 15})
    \end{align*}
  \end{exambox}
\end{frame}

% ===========================================================================
%  SECTION 3: ทศนิยมและร้อยละ
% ===========================================================================
\section{ทศนิยมและร้อยละ (Decimals \& Percentages)}

\begin{frame}{ตัวอย่าง: การแปลง}
  \begin{columns}[T]
    \column{0.33\textwidth}
      \centering
      \textbf{เศษส่วน $\to$ ทศนิยม}
      \[
        \frac{1}{2} = 0.5
      \]
      \[
        \frac{3}{4} = 0.75
      \]
      \[
        \frac{1}{3} = 0.333\ldots
      \]

    \column{0.33\textwidth}
      \centering
      \textbf{ทศนิยม $\to$ ร้อยละ}
      \[
        0.5 \times 100\% = 50\%
      \]
      \[
        0.25 \times 100\% = 25\%
      \]
      \[
        1.5 \times 100\% = 150\%
      \]

    \column{0.33\textwidth}
      \centering
      \textbf{ร้อยละ $\to$ เศษส่วน}
      \[
        50\% = \frac{50}{100} = \frac{1}{2}
      \]
      \[
        25\% = \frac{25}{100} = \frac{1}{4}
      \]
      \[
        10\% = \frac{10}{100} = \frac{1}{10}
      \]
  \end{columns}
\end{frame}

% ===========================================================================
%  SECTION 4: อัตราส่วนและสัดส่วน
% ===========================================================================
\section{อัตราส่วนและสัดส่วน (Ratios \& Proportions)}

\begin{frame}{อัตราส่วน (Ratio)}
  \begin{explanation}[นิยาม]
    \textbf{อัตราส่วน} $a : b$ คือการเปรียบเทียบปริมาณ $a$ ต่อ $b$
    เขียนเป็น $\frac{a}{b}$ ก็ได้ (อ่านว่า ``$a$ ต่อ $b$'')
  \end{explanation}

  \begin{exambox}[ตัวอย่าง]
    \begin{itemize}
      \item ในห้องมีชาย 15 คน หญิง 10 คน
      \par อัตราส่วน ชาย : หญิง = $15 : 10 = 3 : 2$
      \item นํ้าเชื่อมใช้นํ้าตาล 2 ส่วน นํ้า 5 ส่วน
      \par อัตราส่วน นํ้าตาล : นํ้า = $2 : 5$
      \item แผนที่มาตราส่วน $1 : 50{,}000$
      \par 1 ซม. ในแผนที่ = 50,000 ซม. ในความเป็นจริง
    \end{itemize}
  \end{exambox}
\end{frame}

\begin{frame}{สัดส่วน (Proportion)}
  \begin{explanation}[นิยาม]
    \textbf{สัดส่วน} คือการเท่ากันของสองอัตราส่วน: \quad $a : b = c : d$

    \medskip
    เขียนอีกแบบ: $\dfrac{a}{b} = \dfrac{c}{d}$
  \end{explanation}

  \begin{exambox}[ตัวอย่าง: การหาค่าตัวแปรในสัดส่วน]
    จงหา $x$ เมื่อ $3 : 5 = x : 20$

    \medskip
    \textbf{วิธีทำ:} $\dfrac{3}{5} = \dfrac{x}{20}$

    คูณไขว้: $3 \times 20 = 5 \times x$

    $60 = 5x$ \quad $\Rightarrow$ \quad $x = 12$
  \end{exambox}
\end{frame}

\begin{frame}{สัดส่วนตรงและสัดส่วนผกผัน}
  \begin{columns}[T]
    \column{0.48\textwidth}
      \textbf{สัดส่วนตรง (Direct Proportion)}
      \begin{itemize}
        \item $y \propto x$
        \item $y = kx$ ($k$ = ค่าคงที่)
        \item $x$ เพิ่ม $\rightarrow$ $y$ เพิ่ม
        \item เช่น ระยะทางกับเวลา (ความเร็วคงที่)
      \end{itemize}

      \begin{exambox}[ตัวอย่าง]
        ขนม 3 ชิ้น 15 บาท \\
        ขนม 5 ชิ้น = ? บาท \\[0.3em]
        $\frac{3}{15} = \frac{5}{x}$ \\
        $x = 25$ บาท
      \end{exambox}

    \column{0.48\textwidth}
      \textbf{สัดส่วนผกผัน (Inverse Proportion)}
      \begin{itemize}
        \item $y \propto \frac{1}{x}$
        \item $y = \frac{k}{x}$ ($k$ = ค่าคงที่)
        \item $x$ เพิ่ม $\rightarrow$ $y$ ลด
        \item เช่น ความเร็วกับเวลา (ระยะทางคงที่)
      \end{itemize}

      \begin{exambox}[ตัวอย่าง]
        คน 4 คน ทำงานเสร็จใน 6 วัน \\
        คน 8 คน = ? วัน \\[0.3em]
        $4 \times 6 = 8 \times x$ \\
        $x = 3$ วัน
      \end{exambox}
  \end{columns}
\end{frame}

% ===========================================================================
%  SECTION 5: สมบัติของจำนวนจริง
% ===========================================================================
\section{สมบัติของจำนวนจริง (Properties of Real Numbers)}

\begin{frame}{สมบัติการสลับที่และการเปลี่ยนหมู่}
  \begin{explanation}[สมบัติการสลับที่ (Commutative)]
    \[
      a + b = b + a \qquad\qquad a \times b = b \times a
    \]
  \end{explanation}

  \begin{explanation}[สมบัติการเปลี่ยนหมู่ (Associative)]
    \[
      (a + b) + c = a + (b + c) \qquad\qquad (a \times b) \times c = a \times (b \times c)
    \]
  \end{explanation}

  \begin{exambox}[ตัวอย่าง]
    \begin{itemize}
      \item Commutative: $3 + 5 = 5 + 3 = 8$, \quad $4 \times 7 = 7 \times 4 = 28$
      \item Associative: $(2 + 3) + 4 = 2 + (3 + 4) = 9$
    \end{itemize}
  \end{exambox}
\end{frame}

\begin{frame}{สมบัติการแจกแจง (Distributive)}
  \begin{explanation}[สมบัติการแจกแจง]
    \[
      a \times (b + c) = a \times b + a \times c
    \]
    ใช้ได้ทั้งจากซ้ายและขวา: $(b + c) \times a = b \times a + c \times a$
  \end{explanation}

  \begin{exambox}[ตัวอย่าง]
    \begin{align*}
      3 \times (4 + 5) &= 3 \times 4 + 3 \times 5 \\
      3 \times 9 &= 12 + 15 \\
      27 &= 27 \quad \checkmark
    \end{align*}

    \textbf{ใช้กลับด้าน:} $2x + 3x = (2+3)x = 5x$
  \end{exambox}
\end{frame}

\begin{frame}{สมบัติเอกลักษณ์ (Identity)}
  \begin{explanation}[สมบัติเอกลักษณ์]
    \textbf{เอกลักษณ์} คือค่าที่เมื่อนำไปดำเนินการกับจำนวนใด ๆ แล้ว \highlight{ได้จำนวนนั้นกลับคืนมา}
  \end{explanation}

  \begin{columns}[T]
    \column{0.48\textwidth}
      \textbf{เอกลักษณ์การบวก: 0}
      \begin{itemize}
        \item $a + 0 = a$
        \item $0 + a = a$
        \item \textbf{เหตุผล:} การ ``บวก 0'' คือการไม่เพิ่มอะไรเลย
      \end{itemize}

      \begin{exambox}[ตัวอย่าง]
        $7 + 0 = 7$ \\
        $0 + (-3) = -3$ \\
        $x + 0 = x$
      \end{exambox}

    \column{0.48\textwidth}
      \textbf{เอกลักษณ์การคูณ: 1}
      \begin{itemize}
        \item $a \times 1 = a$
        \item $1 \times a = a$
        \item \textbf{เหตุผล:} การ ``คูณ 1'' คือการไม่เปลี่ยนขนาด
      \end{itemize}

      \begin{exambox}[ตัวอย่าง]
        $5 \times 1 = 5$ \\
        $1 \times (-2) = -2$ \\
        $\frac{1}{2} \times 1 = \frac{1}{2}$
      \end{exambox}
  \end{columns}
\end{frame}

\begin{frame}{สมบัติอินเวอร์ส (Inverse)}
  \begin{explanation}[สมบัติอินเวอร์ส]
    \textbf{อินเวอร์ส} คือค่าที่เมื่อนำไปดำเนินการกับจำนวนใด ๆ แล้ว \highlight{ได้เอกลักษณ์กลับคืนมา}
  \end{explanation}

  \begin{columns}[T]
    \column{0.48\textwidth}
      \textbf{อินเวอร์สการบวก: $-a$}
      \begin{itemize}
        \item $a + (-a) = 0$
        \item $(-a) + a = 0$
        \item เรียก $-a$ ว่า \highlight{จำนวนตรงข้าม}
      \end{itemize}

      \begin{exambox}[ตัวอย่าง]
        $5 + (-5) = 0$ \\
        $(-3) + 3 = 0$ \\
        $-\frac{1}{2} + \frac{1}{2} = 0$
      \end{exambox}

    \column{0.48\textwidth}
      \textbf{อินเวอร์สการคูณ: $\frac{1}{a}$}
      \begin{itemize}
        \item $a \times \frac{1}{a} = 1$ \; ($a \neq 0$)
        \item เรียก $\frac{1}{a}$ ว่า \highlight{ส่วนกลับ}
        \item \textbf{0 ไม่มีอินเวอร์สการคูณ!}
      \end{itemize}

      \begin{exambox}[ตัวอย่าง]
        $4 \times \frac{1}{4} = 1$ \\
        $(-2) \times (-\frac{1}{2}) = 1$ \\
        $\frac{2}{3} \times \frac{3}{2} = 1$
      \end{exambox}
  \end{columns}
\end{frame}

\begin{frame}{ตัวอย่าง: ตัวดำเนินการใหม่ $\star$}
  \begin{explanation}[นิยามตัวดำเนินการ $\star$]
    กำหนด $a \star b = a + 2b$ (ตัวอย่างตัวดำเนินการใหม่ที่ไม่ใช่ $+$ หรือ $\times$)
  \end{explanation}

  \begin{columns}[T]
    \column{0.5\textwidth}
      \textbf{ตรวจสอบสมบัติการสลับที่:}
      \begin{align*}
        a \star b &= a + 2b \\
        b \star a &= b + 2a
      \end{align*}
      $a + 2b \neq b + 2a$ 

      \medskip
      \highlight{$\star$ ไม่มีสมบัติการสลับที่}

    \column{0.5\textwidth}
      \textbf{ตรวจสอบสมบัติการเปลี่ยนหมู่:}
      \begin{align*}
        (a \star b) \star c &= (a+2b) \star c \\
        &= a + 2b + 2c \\
        a \star (b \star c) &= a \star (b+2c) \\
        &= a + 2b + 4c
      \end{align*}
      แตกต่างกัน!

      \medskip
      \highlight{$\star$ ไม่มีสมบัติการเปลี่ยนหมู่}
  \end{columns}
\end{frame}

\begin{frame}{ตัวอย่าง: ตัวดำเนินการ $\star$ (ต่อ)}
  \begin{exambox}[ตรวจสอบสมบัติการแจกแจง]
    \textbf{แจกแจงทางซ้าย:} $a \star (b + c) = a + 2(b+c) = a + 2b + 2c$

    $(a \star b) + (a \star c) = (a+2b) + (a+2c) = 2a + 2b + 2c$

    \medskip
    $a \star (b+c) \neq (a \star b) + (a \star c)$ โดยทั่วไป

    \highlight{$\star$ ไม่มีสมบัติการแจกแจง}
  \end{exambox}

  \begin{intuition}[บทเรียนจาก $\star$]
    \begin{itemize}
      \item ตัวดำเนินการที่เราคุ้นเคย ($+, \times$) มีสมบัติมากมาย
      \item แต่ตัวดำเนินการใหม่ที่เรานิยามขึ้นเอง \textbf{อาจไม่มีสมบัติเหล่านี้}
      \item ต้อง \highlight{ตรวจสอบทีละสมบัติ} — อย่าคิดว่าเป็นจริงเสมอไป!
    \end{itemize}
  \end{intuition}
\end{frame}

\begin{frame}{ตรวจสอบ: $\star$ มีเอกลักษณ์และอินเวอร์สไหม?}
  \begin{exambox}[เอกลักษณ์ของ $\star$]
    \textbf{เอกลักษณ์การบวก:} $a \star e = a$ หมายถึง $a + 2e = a \Rightarrow 2e = 0 \Rightarrow e = 0$

    $a \star 0 = a + 2(0) = a$ \quad $\checkmark$

    แต่ $0 \star a = 0 + 2a \neq a$ \quad $\times$ \;(ไม่สลับที่!)

    \medskip
    \highlight{$\star$ ไม่มีเอกลักษณ์ที่แท้จริง} เพราะ $\star$ ไม่มีสมบัติการสลับที่
  \end{exambox}

  \begin{exambox}[อินเวอร์สของ $\star$]
    ไม่สามารถหาอินเวอร์สได้ ถ้าไม่มีเอกลักษณ์ที่แน่นอน

    \medskip
    \highlight{บทเรียน:} สมบัติทั้งหลายเชื่อมโยงกัน —
    ถ้าไม่มีสมบัติการสลับที่ ก็ส่งผลต่อสมบัติอื่น ๆ ด้วย
  \end{exambox}
\end{frame}

% ===========================================================================
%  SECTION 6: ค่าสัมบูรณ์
% ===========================================================================
\section{ค่าสัมบูรณ์ (Absolute Value)}

\begin{frame}{ค่าสัมบูรณ์คืออะไร?}
  \begin{explanation}[นิยาม]
    \textbf{ค่าสัมบูรณ์} ของ $x$ ($|x|$) คือ \highlight{ระยะห่างจาก 0} บนเส้นจำนวน

    \[
      |x| =
      \begin{cases}
        x,  & \text{ถ้า } x \geq 0 \\
        -x, & \text{ถ้า } x < 0
      \end{cases}
    \]
  \end{explanation}

  \begin{exambox}[ตัวอย่าง]
    $|5| = 5$ \qquad $|-3| = 3$ \qquad $|0| = 0\qquad |2 - 7| = |-5| = 5$ \quad (ระยะห่างของ 2 และ 7)
  \end{exambox}

  \begin{intuition}[คิดง่าย ๆ]
    ค่าสัมบูรณ์ = \highlight{ค่าบวกของจำนวนนั้น} (หรือศูนย์)
    — มันคือการถามว่า ``ห่างจาก 0 เท่าไหร่?''
  \end{intuition}
\end{frame}

\begin{frame}{สมบัติของค่าสัมบูรณ์}
  \begin{explanation}[สมบัติที่สำคัญ]
    \begin{enumerate}
      \item $|a| \geq 0$ \quad (ค่าสัมบูรณ์ไม่เป็นลบเสมอ)
      \item $|a| = |-a|$
      \item $|ab| = |a| \cdot |b|$
      \item $|a + b| \leq |a| + |b|$ \quad (อสมการสามเหลี่ยม — Triangle Inequality)
      \item $\sqrt{a^2} = |a|$ \quad (\textbf{ระวัง!} ไม่ใช่ $a$ เฉย ๆ)
    \end{enumerate}
  \end{explanation}

  \begin{exambox}[ตัวอย่างการใช้สมบัติ]
    \begin{itemize}
      \item $|3 \times (-4)| = |3| \times |-4| = 3 \times 4 = 12$
      \item $|3 + (-5)| = |-2| = 2 \leq |3| + |-5| = 8$ \quad $\checkmark$
    \end{itemize}
  \end{exambox}
\end{frame}

\begin{frame}{การแก้สมการค่าสัมบูรณ์}
  \begin{explanation}[หลักการแก้สมการ]
    $|x| = a$ ($a \geq 0$)
    \[
      |x| = a \quad \iff \quad x = a \;\; \text{หรือ} \;\; x = -a
    \]
    \textbf{หมายเหตุ:} ถ้า $a < 0$ แสดงว่า \textbf{ไม่มีคำตอบ} เพราะค่าสัมบูรณ์ $\geq 0$ เสมอ
  \end{explanation}
\end{frame}

\begin{frame}{ตัวอย่าง: การแก้สมการค่าสัมบูรณ์}
  \begin{exambox}[โจทย์ 1]
    $|x - 3| = 5$
    \[
      x - 3 = 5 \quad \text{หรือ} \quad x - 3 = -5
    \]
    \[
      x = 8 \quad \text{หรือ} \quad x = -2
    \]
  \end{exambox}

  \begin{exambox}[โจทย์ 2]
    $|2x + 1| = 7$
    \[
      2x + 1 = 7 \;\; \Rightarrow \;\; 2x = 6 \;\; \Rightarrow \;\; x = 3
    \]
    \[
      2x + 1 = -7 \;\; \Rightarrow \;\; 2x = -8 \;\; \Rightarrow \;\; x = -4
    \]
    \textbf{คำตอบ:} $x = 3$ หรือ $x = -4$
  \end{exambox}
\end{frame}

\begin{frame}{การแก้อสมการค่าสัมบูรณ์}
  \begin{explanation}[หลักการแก้อสมการ]
    \textbf{แบบที่ 1:} $|x| < a$ \quad ($a > 0$)
    \[
      -a < x < a
    \]
    (ระยะห่างจาก 0 น้อยกว่า $a$)

    \textbf{แบบที่ 2:} $|x| > a$ \quad ($a > 0$)
    \[
      x < -a \quad \textbf{หรือ} \quad x > a
    \]
    (ระยะห่างจาก 0 มากกว่า $a$)
  \end{explanation}
\end{frame}

\begin{frame}{ตัวอย่าง: อสมการ $|x - 2| \leq 4$}
  \begin{exambox}[โจทย์: $|x - 2| \leq 4$]
    \[
      -4 \leq x - 2 \leq 4
    \]
    \[
      -2 \leq x \leq 6
    \]
    \textbf{คำตอบ:} $x \in [-2, 6]$
  \end{exambox}

  \begin{intuition}[ความหมาย]
    $x$ ห่างจาก 2 ไม่เกิน 4 หน่วย $\rightarrow$ ระยะจาก 2 ถึง -2 คือ 4, ระยะจาก 2 ถึง 6 คือ 4
  \end{intuition}
\end{frame}

\begin{frame}{ตัวอย่าง: อสมการ $|3x + 1| > 5$}
  \begin{exambox}[โจทย์: $|3x + 1| > 5$]
    \[
      3x + 1 < -5 \quad \text{หรือ} \quad 3x + 1 > 5
    \]
    \[
      x < -2 \quad \text{หรือ} \quad x > \frac{4}{3}
    \]
    \textbf{คำตอบ:} $x \in (-\infty, -2) \cup (\frac{4}{3}, \infty)$
  \end{exambox}
\end{frame}

\begin{frame}{แนวทางการแก้โจทย์ค่าสัมบูรณ์}
  \begin{intuition}[ขั้นตอนการแก้โจทย์ค่าสัมบูรณ์]
    \begin{enumerate}
      \item \textbf{จัดรูป}ให้อยู่ในรูป $|\text{นิพจน์}| = a$, $|\text{นิพจน์}| < a$, หรือ $|\text{นิพจน์}| > a$
      \item \textbf{แยกกรณี}ตามนิยาม:
      \begin{itemize}
        \item $|X| = a \;\rightarrow\; X = a$ หรือ $X = -a$
        \item $|X| < a \;\rightarrow\; -a < X < a$
        \item $|X| > a \;\rightarrow\; X < -a$ หรือ $X > a$
      \end{itemize}
      \item \textbf{แก้}สมการ/อสมการในแต่ละกรณี และ \textbf{ตรวจคำตอบ}
    \end{enumerate}
  \end{intuition}
\end{frame}

\begin{frame}{ตัวอย่าง: ข้อสอบจริง}
  \begin{exambox}[ข้อสอบจริงปี 68 ข้อ 14]
    อะไรคือคำตอบของสมการ $\left\lvert 2x-3\right\rvert = 5x+6$
  \end{exambox}
\end{frame}

\begin{frame}{ตัวอย่าง: ข้อสอบจริง}
  \begin{solbox}
    \begin{align*}
      |2x-3| &= 5x+6 \quad \text{โดยมีเงื่อนไขว่า } 5x+6 \geq 0 \implies x \geq -\frac{6}{5} \\
      \textbf{กรณีที่ 1:} \quad 2x-3 &= 5x+6 \\
      -3x &= 9 \\
      x &= -3 \quad \text{(ใช้ไม่ได้เพราะขัดกับเงื่อนไข } x \geq -1.2\text{)} \\
      \textbf{กรณีที่ 2:} \quad -(2x-3) &= 5x+6 \\
      -2x+3 &= 5x+6 \\
      -7x &= 3 \\
      x &= -\frac{3}{7} \quad \text{(ใช้ได้เพราะสอดคล้องกับเงื่อนไข } x \geq -1.2\text{)} \\
    \end{align*}
  \end{solbox}
\end{frame}
% ===========================================================================
%  SECTION 7: เลขยกกำลังและราก
% ===========================================================================
\section{เลขยกกำลังและราก (Exponents \& Radicals)}

\begin{frame}{เลขยกกำลัง (Exponents)}
  \begin{explanation}[นิยาม]
    $a^n$ = $a$ คูณกัน $n$ ครั้ง \quad ($n \in \N$)

    \medskip
    \begin{itemize}
      \item $a^1 = a$
      \item $a^0 = 1$ \quad ($a \neq 0$)
      \item $a^{-n} = \dfrac{1}{a^n}$ \quad ($a \neq 0$)
    \end{itemize}
  \end{explanation}

  \begin{exambox}[ตัวอย่าง]
    $2^3 = 2 \times 2 \times 2 = 8$ \qquad
    $5^0 = 1$

    $10^{-1} = \frac{1}{10} = 0.1$ \qquad
    $2^{-3} = \frac{1}{2^3} = \frac{1}{8}$
  \end{exambox}
\end{frame}

\begin{frame}{สมบัติของเลขยกกำลัง}
  \begin{explanation}[กฎของเลขยกกำลัง]
    \[
    \begin{array}{c|c}
      \textbf{กฎ} & \textbf{ตัวอย่าง} \\
      \hline
      a^m \cdot a^n = a^{m+n}         & 2^3 \cdot 2^4 = 2^7 = 128 \\[0.4em]
      \dfrac{a^m}{a^n} = a^{m-n}     & \dfrac{2^5}{2^2} = 2^3 = 8 \\[0.4em]
      (a^m)^n = a^{mn}               & (2^3)^2 = 2^6 = 64 \\[0.4em]
      (ab)^n = a^n \cdot b^n         & (2 \times 3)^2 = 2^2 \times 3^2 = 4 \times 9 = 36 \\[0.4em]
      \left(\dfrac{a}{b}\right)^n = \dfrac{a^n}{b^n} & \left(\dfrac{2}{3}\right)^2 = \dfrac{4}{9}
    \end{array}
    \]
  \end{explanation}
\end{frame}

\begin{frame}{ราก (Radicals)}
  \begin{explanation}[นิยาม]
    $\sqrt[n]{a}$ คือจำนวนที่ยกกำลัง $n$ แล้วได้ $a$

    \medskip
    $\sqrt{a}$ = $\sqrt[2]{a}$ (รากที่สอง) \quad — \quad $(\sqrt{a})^2 = a$ \; ($a \geq 0$)
  \end{explanation}

  \begin{exambox}[ตัวอย่าง]
    $\sqrt{9} = 3$ \quad เพราะ $3^2 = 9$

    $\sqrt[3]{8} = 2$ \quad เพราะ $2^3 = 8$

    $\sqrt{16} = 4$ \quad (ไม่ใช่ $\pm 4$ — รากที่สองที่เป็นบวก)

    \medskip
    \textbf{ระวัง!} $\sqrt{x^2} = |x|$ ไม่ใช่ $x$ เฉย ๆ
  \end{exambox}
\end{frame}

\begin{frame}{สมบัติของราก}
  \begin{explanation}[กฎของราก]
    \[
    \begin{array}{c|c}
      \textbf{กฎ} & \textbf{ตัวอย่าง} \\
      \hline
      \sqrt[n]{ab} = \sqrt[n]{a} \cdot \sqrt[n]{b} & \sqrt{12} = \sqrt{4 \times 3} = 2\sqrt{3} \\[0.4em]
      \sqrt[n]{\dfrac{a}{b}} = \dfrac{\sqrt[n]{a}}{\sqrt[n]{b}} & \sqrt{\dfrac{4}{9}} = \dfrac{2}{3} \\[0.4em]
      \sqrt[m]{\sqrt[n]{a}} = \sqrt[mn]{a} & \sqrt{\sqrt{16}} = \sqrt[4]{16} = 2
    \end{array}
    \]
  \end{explanation}

  \begin{exambox}[การเขียนรากเป็นเลขยกกำลัง]
    $\sqrt[n]{a} = a^{1/n}$ \qquad $\sqrt[n]{a^m} = a^{m/n}$

    \medskip
    ตัวอย่าง: $\sqrt[3]{x^2} = x^{2/3}$ \qquad $\sqrt{x} = x^{1/2}$
  \end{exambox}
\end{frame}

% ===========================================================================
%  SUMMARY
% ===========================================================================
\section{สรุป}

\begin{frame}{สรุปเนื้อหา}
  \begin{explanation}[สิ่งที่ได้เรียนรู้]
    \begin{enumerate}
      \item \textbf{ระบบจำนวน}: $\N \subset \Z \subset \Q \subset \R$ — ตรรกยะ vs อตรรกยะ
      \item \textbf{เศษส่วน}: การบวก ลบ คูณ หาร และการทำให้เป็นอย่างตํ่า
      \item \textbf{ร้อยละ \& สัดส่วน}: การแปลงระหว่างเศษส่วน-ทศนิยม-ร้อยละ, สัดส่วนตรงและผกผัน
      \item \textbf{สมบัติของ $\R$}: Commutative, Associative, Distributive, Identity, Inverse
      \item \textbf{ตัวดำเนินการใหม่}: ต้องตรวจสอบสมบัติทีละข้อ — อย่าคิดว่าเป็นจริงเสมอ
      \item \textbf{ค่าสัมบูรณ์}: $|x| = $ ระยะห่างจาก 0, การแก้สมการและอสมการ
      \item \textbf{เลขยกกำลัง \& ราก}: $a^m \cdot a^n = a^{m+n}$, $\sqrt{x^2} = |x|$
    \end{enumerate}
  \end{explanation}

  \begin{center}
    \large เอกสารนี้เป็นส่วนหนึ่งของเนื้อหาเตรียมสอบ \textbf{สอวน. คอมพิวเตอร์}
  \end{center}
\end{frame}

\end{document}
